\section{Source reconciliation and a further research programme}
\label{sec:audit}

\subsection{What the audit establishes}

The source claims used in the proofs were checked against the pinned
revision, including the finite-cutoff conventions and the exact matrices.
The audit establishes one stale research-status statement and a precise
centering reconciliation. It did not establish an additional false theorem
in the inspected parts of the six incoming reports. This conclusion is
limited to those parts: acquiring the full source inventory does not
constitute checking every proof in the collective manuscript.

\begin{longtable}{@{}>{\raggedright\arraybackslash}p{0.26\textwidth}>{\raggedright\arraybackslash}p{0.38\textwidth}>{\raggedright\arraybackslash}p{0.28\textwidth}@{}}
\toprule
Source & Material checked for this continuation & Result \\
\midrule
\endfirsthead
\toprule
Source & Material checked for this continuation & Result \\
\midrule
\endhead
Canonical fifth-moment report \cite{Canonical} &
Tail blocks, inverse transform, diagonal rank, and the first-$q$ conjecture &
Normalization retained; gap four advances its conjecture \\
Differential polygamma and Lerch jets \cite{Differential} &
Adjacent-order proof and fixed-gap recurrence &
Recurrence and determinant signs retained \\
Bernoulli--Tornheim--digamma \cite{Bernoulli} &
Fixed smaller orders $2,4,6$; shifted-Hankel reduction; tail foundation;
Tornheim subtraction; differential and research sections &
No additional false formula found; R8 status needs reconciliation \\
Polygamma tail identities \cite{Polygamma} &
Moment rank and reconstruction; spectral polynomial formulas; sixth
mixed moment &
Supplies the already available spectral coordinates used below \\
Cubic tail descent \cite{Cubic} &
Cubic Gamma product and the master identity with its subtraction &
No error established in the inspected product and limit arguments \\
Determinantal harmonic resonance \cite{Resonance} &
First resonance and the cutoff-to-spectral transfer law &
The amplitude-derivative correction remains necessary \\
Multishift reflection \cite{Multishift} &
Entire-order transform and full horizontal-cut reflection &
No error established in the inspected holomorphy and jump conventions \\
\bottomrule
\end{longtable}

The accompanying source-audit record identifies the exact section files
and equation labels. The preserved Canonical report is already under
\path{Analysis/Polylogarithms/docs/reports/stieltjes-harmonic-resonance/continuations/canonical-digamma-fifth/};
it is not one of the six current incoming archives. The earlier extracted
tail-rank source used during this work was checked against that preserved
file and has the identical Git blob.

\subsection{A sixth mixed moment whose spectral evaluation already exists}
\label{sh:audit-sixth}

Research item R8 of the Bernoulli--Tornheim--digamma report asks for a
harmonic-spectral evaluation of
\[
 \widetilde M_6=\FP\int_0^1(\psi(x)+\gamma)^2\psi'(x)^2\,dx.
\]
Here and throughout this subsection, $\FP\int_0^1f(x)\,dx$ means the
coefficient of $\varepsilon^0(\log\varepsilon)^0$ in
$\int_\varepsilon^1f(x)\,dx$ as $\varepsilon\downarrow0$, in the
original coordinate $x$. The operation is linear on the Laurent-log
expansions used here. The existence question is already answered by
the Polygamma Tail Identities report \cite{Polygamma}. Its equation
\texttt{sixth:I} evaluates the uncentered mixed integral; equations
\texttt{spectral:exampleDP} and
\texttt{spectral:digamma\_trigammasquare} evaluate the lower mixed term.
The differing centering can be reconciled explicitly.

Define, in decreasing-index convention,
\begin{align}\label{audit:DP}
 D_P(s)={}&-4\zeta(3,s+1)-4\zeta(s+4)
          +2\zeta(s+2,2)+2\zeta(2)\zeta(s+2)\notag\\
         &+4\zeta(s+3,1)-4\gamma\zeta(s+3),\qquad \Re s>-1.
\end{align}
All the displayed series are holomorphic in this half-plane. To specify
the other inherited coordinate without leaving an unnamed integral,
put, for $n\ge0$,
\begin{align*}
 b_n&=H_n-\gamma,&c_n&=H_n^{(2)}+\zeta(2),&
 d_n&=H_n^{(3)}-\zeta(3),\\
 e_n&=H_n^{(4)}+\zeta(4),&f_n&=H_n^{(5)}-\zeta(5),
\end{align*}
and
\begin{align}
 r_n={}&4b_n^2d_n+2b_nc_n^2-6c_nd_n-10b_ne_n+6f_n,
       \label{audit:residue}\\
 \Omega_5={}&-\sum_{n\ge1}(r_n-r_{n-1})\log n
           =\sum_{n\ge1}r_n\log(1+1/n).\label{audit:Omega}
\end{align}
These are ordinary absolutely convergent series. For example, the
harmonic-tail expansions and $8\zeta(2)^2=20\zeta(4)$ give
$r_n=O(\log(n+1)/n+\log^2(n+1)/n^2)$ and
$r_n-r_{n-1}=O(\log(n+1)/n^2)$; summation by parts has a vanishing
boundary term. These definitions agree exactly with the source's
\texttt{sixth:D} and \texttt{sixth:omega-series}.

\begin{proposition}[Reconciliation of the centered sixth coordinate]
\label{audit:sixth}
The inherited spectral evaluations imply
\begin{equation}\label{sh:centered-sixth-reconciliation}
 \boxed{\widetilde M_6
   =\Omega_5+2\gamma D_P'(0)+4\gamma^2\zeta'(3)
      +\frac{31}{3}\zeta(2)\zeta(3)-\frac{19}{3}\zeta(5).}
\end{equation}
\end{proposition}
\begin{proof}
The three source evaluations in the same finite-part convention are
\begin{align*}
 \FP\int_0^1\psi^2\psi'^2\,dx
 &=\Omega_5+2\gamma^2\zeta(3)-\frac{35}{3}\gamma\zeta(4)
       +\frac{31}{3}\zeta(2)\zeta(3)-\frac{19}{3}\zeta(5),\\
 \FP\int_0^1\psi\psi'^2\,dx
 &=\frac{35}{6}\zeta(4)-2\gamma\zeta(3)+D_P'(0),\\
 \FP\int_0^1\psi'^2\,dx
 &=2\zeta(3)+4\zeta'(3).
\end{align*}
Expand $(\psi+\gamma)^2\psi'^2$ and use linearity. The coefficients
of $\gamma\zeta(4)$ and $\gamma^2\zeta(3)$ cancel exactly, leaving
\eqref{sh:centered-sixth-reconciliation}. In the Canonical fifth-moment
notation one also has $D_P(s)=2J(s+1)$, by the stuffle product; hence
$D_P'(0)=2J'(1)$.
\end{proof}

The proposed correction to R8 is therefore: determine whether the
specified combination in \eqref{sh:centered-sixth-reconciliation}
reduces to a smaller, explicitly stated collection of constants, such
as selected ordinary zeta derivatives and prescribed height-one
spectral data. Existence of a harmonic-spectral representation should
be marked resolved. This reconciliation does not prove a reduction
to that smaller collection, and it is not an arithmetic-independence
statement about $\Omega_5$.

\subsection{Conventions that must survive integration}

Three distinctions materially affect the values of the formulas.
First, the original centered lattice has
$\zeta(-2j,1/2)=0$, whereas a shifted lattice generally does not.
Theorems~\ref{sh:transform} and \ref{sh:jets} give the two required
corrections explicitly. Second, equality of zero-order continued and
ordinary values does not imply equality after spectral differentiation:
the nonzero $\zeta^{(j)}(-k,\alpha)$ terms must be retained. Third,
$\mathcal K_{+,+}$ and its $(2,2)$ coefficient have an ordinary
conditional interpretation; their original series must not be
described as absolutely convergent. The other sign pairings are
absolutely convergent.

These are normalization requirements for using the new results, not
claims that the inspected predecessors made all three mistakes.
The separately unresolved Gaussian $S_6$ and revised $S_8$ candidates
remain conjectural in this continuation. A proved reduction for a
different tail family supplies no missing identity for those values.

\subsection{Further research questions with concrete targets}
\label{sec:questions}

The following questions concern exact evaluations, structure, or proof
certificates. The numerical evidence accompanying the sources can help
choose cases, but each proposed endpoint is an identity or a rigorous
nonvanishing theorem.

\begin{enumerate}[label=\textbf{Q\arabic*.},ref=Q\arabic*,itemsep=10pt]
\item \textbf{Complete the optimal-prefix conjecture beyond gap four.}
For even $8\le p<q$ with $q-p\ge6$, prove or refute
$\det M_{p,q}\ne0$. The first unproved fixed-gap family is
$(p,q)=(2b-6,2b)$ with $b\ge7$. Its boundary determinant has size
three. The explicit sign change already seen between $b=6$ and $b=7$
rules out a proof based on one fixed sign for this family. A useful
target is an exact expression whose factors prove nonvanishing while
allowing the determinant's sign to vary. Comparing the size-$c$
fixed-gap determinant with the size-$p/2$ complementary determinant
in the shifted-Hankel approach may reduce different parameter ranges.

\item \textbf{Classify higher minors of the fundamental recurrence.}
Theorem~\ref{eg:two-minors} is uniform in the gap and the two chosen
columns. Determine which larger minors vanish or change sign, and
whether a recurrence for their exterior products separates those two
possibilities. The one-negative-index induction in the present proof
does not directly extend to several sign changes. A successful
replacement should explain the recorded third-order sign reversal,
not impose a positivity property contradicted by it.

\item \textbf{Determine the exact formal rank after a Hurwitz shift.}
Section~\ref{sh:section} gives at most $q+2$ coordinates over
$\Q(v)$. Is that list generically minimal, and how many initial
even moments reconstruct it? Describe every exceptional shift in a
stated real or complex domain, including the known degeneration at
$v=1$. This is a question about rational coefficient matrices; it
does not require independence of the specialized Hurwitz values.
The Bernoulli-polynomial endpoint and the highest single-Hurwitz
column should be kept throughout any determinant computation.

\item \textbf{Allow different shifts in the two factors.}
Replace the product by
$\zeta(p,n+v)\zeta(q,n+w)$. Integer values of $v-w$ can be
related by finite tail corrections. For general rational differences,
determine an explicit finite normal form and the necessary depth-two
shifted coordinates. The first test is whether a discrete primitive
analogous to Theorem~\ref{sh:primitive} closes in single Hurwitz
values and endpoint products, and exactly which shift differences
permit that closure.

\item \textbf{Combine tail-order derivatives with the weight jets.}
Theorem~\ref{sh:jets} differentiates the factor $x_n^{-s}$.
Derivatives in $p$ or $q$ instead create logarithmic harmonic tails.
Construct the mixed Taylor-Laurent expansion in all three variables,
including motion of the polar hyperplanes and coincidences of
resonances. The target is a finite expression in specified derivatives
of single and multiple Hurwitz functions and generalized Stieltjes
constants, with the scale and amplitude corrections visible at every
order. This would extend the source's Lerch and Tornheim jet machinery
to the shifted product setting.

\item \textbf{Replace signs by higher roots of unity.}
For a fixed conductor $N$, study the analogues of
$\mathcal K_{\alpha,\beta}$ with $N$th roots of unity as colors
and a corresponding outer character. State explicitly which
conjugation and weight-parity sectors reduce in depth. The level-two
parity formula used here cannot simply be reused unchanged. A
concrete objective is a conductor-four generator expressed through
digamma at quarter shifts and a specified collection of cyclotomic
polylogarithms. Such a result would give an exact bridge toward the
Gaussian setting without assuming the outstanding Gaussian identities.

\item \textbf{Find compact normal forms for differentiated resolvents.}
Parameter differentiation of Theorem~\ref{alt:thm:master} already
gives sums of products of arbitrary repeated rational resolvents.
Can their right sides be organized by finite Bell-polynomial formulas
that display the removable singularities before specialization?
At rational parameters, seek explicit formulas in algebraic
coefficients, logarithms, and Dirichlet $L$-values arising from
polygamma at rational arguments. A separate useful target is a direct
contour or hypergeometric proof of the diagonal square and its Gamma
primitive, independent of the general parity theorem.

\item \textbf{Sum the higher odd weights into one generating identity.}
The finite reductions in \eqref{alt:eq:weighted-tail} handle every
admissible odd polynomial weight. Introduce a third variable for
that degree and determine a compact generating function, or an
explicit differential operator on $O_\pm$. The admissible degree
condition must be built into the formulation: extending it crosses
new convergence thresholds and requires additional subtractions.
An exact endpoint prescription is part of the desired identity.

\item \textbf{Extend the discrete primitive to higher products.}
For three or more Hurwitz tails, finite differences introduce products
of fewer tails. Determine which polynomial weights allow this
hierarchy to terminate in single Hurwitz values, and which force
genuine multiple-Hurwitz coordinates. Hoffman's higher-product
decomposition \cite{Hoffman} supplies the ordered-sum framework.
A useful initial problem is a complete classification for three
equal even orders and weights of degree below the first resonance.

\item \textbf{Reduce the reconciled sixth spectral coordinate.}
Choose a target collection first, then test whether
$\Omega_5+2\gamma D_P'(0)+4\gamma^2\zeta'(3)$ belongs to its
rational span. Candidate collections can include specified ordinary
zeta derivatives and the height-one slopes already defined by the
Gamma bridge. The source gives convergent series and a finite
multiple-zeta derivative representation for the current coordinate;
the missing step is a proved functional reduction among those
representations. Numeric relation searches alone cannot decide it.

\item \textbf{Connect the colored primitive to endpoint periods.}
Proposition~\ref{alt:prop:GF} fixes both an analytic function and its
additive constant in the disk. Continue it to selected boundary
points using a stated path, compute its endpoint limits, and compare
the resulting depth-three polylogarithms under shuffle and stuffle
relations. Determine which combinations admit a Gamma or polygamma
description like the diagonal primitive. At a branch point, the
desired result should include the full finite-part or limiting rule.

\item \textbf{Produce exact certificates for the remaining numerical candidates.}
For the Gaussian $S_6$ and revised $S_8$ conjectures, isolate a finite
algebraic, contour, or hypergeometric identity whose specialization
would prove the claimed value. Separate the convergence argument
from the period reduction, and state the intermediate functions with
fixed branches. The finite telescoper in \eqref{eg:V-difference}
illustrates the desired form of a checkable certificate for a sum;
the Gaussian questions need their own certificates.
\end{enumerate}

The immediate priorities are Q1, Q3, and Q6. They extend the three
proved mechanisms of this article while preserving an exact target
and a clear test for what would constitute a solution.
